\section{Introduction}

We prove existence, uniqueness, and stability, on every fixed finite horizon $T$ and
with no smallness restriction, for the fully coupled
McKean--Vlasov forward--backward system with jumps
\begin{equation}\label{eq:mvfbsdej}
\PaperSystemBox{\left\{\;\begin{aligned}
X_t&=\xi_0+\int_0^t b(s,\Theta_s,\mu_s)\,\diff s+\int_0^t\sigma(s,\Theta_s,\mu_s)\,\diff W_s
  +\int_0^t\!\!\int_E h(s,\Theta_{s}^-,e,\mu_s^-)\,\Nt(\diff s,\diff e),\\[9pt]
Y_t&=g(X_T,\Law(X_T))+\int_t^T f(s,\Theta_s,\mu_s)\,\diff s
  -\int_t^T Z_s\,\diff W_s-\int_t^T\!\!\int_E U_s(e)\,\Nt(\diff s,\diff e),
\end{aligned}\right.}
\end{equation}
driven by a Brownian motion $W$ and an independent compensated Poisson random measure
$\Nt$. The intensity $\nu$ is assumed only $\sigma$-finite: infinite jump
activity, $\nu(E)=\infty$, is admitted. The state space is the
Hilbert space
\[
\mathcal H:=\R^n\times\R^m\times\R^{m\times d}\times L^2(\nu;\R^m),
\]
carrying the norm
\[
|\theta|_{\mathcal H}^2:=|x|^2+|y|^2+|z|_F^2+\Lnorm{u}^2,\qquad \theta=(x,y,z,u),
\]
where $|\cdot|$ is the Euclidean norm (on any $\R^j$) and the Frobenius and
$L^2(\nu;\R^m)$ norms are
\[
|z|_F^2:=\sum_{i,j}z_{ij}^2,
\qquad
\Lnorm{u}^2:=\int_E|u(e)|^2\,\nu(\diff e);
\]
$L^2(\nu;\R^m)$ is henceforth abbreviated $\Lnu$.

We seek a solution tuple
\[
\Theta_t=(X_t,Y_t,Z_t,U_t);
\]
the jump coefficient is evaluated at the predictable version
$\Theta_t^-:=(X_{t-},Y_{t-},Z_t,U_t)$, which takes left limits in the forward and backward
states ($Z,U$ are already predictable). The coefficients are evaluated along the joint-law flows
\[
\mu_t=\Law(\Theta_t),\qquad \mu_t^-:=\Law(\Theta_t^-)\in\mathcal P_2(\mathcal H).
\]
The distinguishing feature is the law argument:
the coefficients are Lipschitz in the quadratic Wasserstein distance $W_2$
on $\mathcal P_2(\mathcal H)$, so the law of the $\Lnu$-valued jump integrand $U$ itself
may enter the coefficients; both the Lipschitz and the monotonicity hypotheses are
imposed only along the solution's own tuple--law pairs $(\Theta,\Law(\Theta))$.

This class
of equations is motivated by the optimality systems of mean-field games and mean-field
control problems with jumps, where the law argument is typically the state or control law,
and by their numerical analysis: the a~posteriori error estimates of Reisinger, Stockinger
and Zhang~\cite{RSZ} for Brownian MV-FBSDEs rest on the Brownian case of the monotonicity
hypothesis used below, and carrying them to jumps needs the well-posedness theory proved
here.

The primary comparison is the theorem of Li and Min~\cite{LiMin}: global well-posedness
for fully coupled mean-field FBSDEs with jumps, proved by continuation under a
jump-extended $G$-monotonicity condition~\cite[(H3.2)]{LiMin}. Their mean-field
interaction is of expectation-functional (independent-copy) type, their jump integrand
enters only through its pointwise values, and their monotonicity hypothesis imposes
finite total jump mass; the present theorem removes all three restrictions, allowing
$W_2$-dependence on the full tuple law, the law of $U$ included, at arbitrary
$\sigma$-finite activity, under the expected $G$-monotonicity condition set out
below. Example~\ref{ex:expfunc} details the embedding of their interaction class.

For Brownian noise, well-posedness of fully coupled mean-field FBSDEs under joint-law
dependence has been studied by Bensoussan, Yam and Zhang~\cite{BYZ}, Chen and
Nie~\cite{ChenNie} and
others~\cite{CarmonaDelarueFBSDE,ChenNieWang,HuaLuo,MinPengQin,TianYu,WuHao,ZhangHuang}
under various monotonicity and weak-coupling assumptions. All of these are Brownian: no
jump channel, no integrand $U$.

With jumps but no mean-field coupling, fully coupled FBSDEs are well posed over arbitrary
horizons by the same monotone-continuation method: Wu~\cite{WuFBSDEP} for the
Brownian--Poisson case, and domination-monotonicity continuation carried to the
L\'evy-driven case by Xu, Tang and Meng~\cite{XuTangMeng}. For jump \emph{mean-field}
systems in decoupled form, the well-posedness and regularity theory is due to
Li~\cite{LiSPA18}. Neither setting addresses the fully coupled problem with full-tuple
law dependence.

Further fully coupled jump mean-field results trade scope for other restrictions.
Matoussi, Manai and Salhi~\cite{MMS} treat $W_2$-dependence on the $(X,Y)$-marginal law
at finite activity, under monotonicity hypotheses and explicit size restrictions on the
Lipschitz constants.
Chen, Shi and Wu~\cite{ChenShiWu} work on an arbitrary fixed horizon in the weakly
coupled regime, where quantitative weak coupling replaces the structural dissipativity
used here. Liu and Zhang~\cite{LiuZhang} allow pointwise dependence on $Y,Z$ but take the
measure argument to be $\Law(X)$, on a small time interval and under finite jump
activity. In a control direction, mean-field-type control driven by jump-diffusions is
solved globally in time by Bensoussan, Huang, Tang and Yam~\cite{BHTY}, with state-law
dependence and finite total jump mass.

Assumption~\ref{ass:mono} below adjoins a jump term to the monotonicity hypothesis
of~\cite[(H.1)(1)]{RSZ}, posed there, as in the condition~\cite[(A1)]{BYZ}, on the joint
law of $(X,Y,Z)$ with a rectangular full-rank $G$. The measure argument is correspondingly the joint law of the full tuple
$(X,Y,Z,U)$, so the law of the $\Lnu$-valued jump integrand may enter the coefficients, and
$\nu$ is $\sigma$-finite, infinite activity included.

The works surveyed above are each restricted
on at least one of the three axes (measure argument, jump activity, coupling strength). We
prove well-posedness under the stated $L^2$ and monotonicity assumptions by global-in-time
monotone continuation; concretely:
\begin{enumerate}[label=(\roman*)]
\item an a~priori stability estimate: the full difference of two solutions is
  controlled by the differences of their coefficients and data (Theorem~\ref{thm:stab});
\item uniqueness in the natural class (Theorem~\ref{thm:uniq});
\item global existence by monotone continuation from a small-coupling base case
  (Theorem~\ref{thm:cwp});
\item robustness: well-posedness survives full-tuple-law perturbations below the
  dissipativity margin at equal forward and backward dimension
  (Corollary~\ref{cor:robust}; Remark~\ref{rem:generalrank} for general rank); and
\item strictness: expected diagonal monotonicity is strictly weaker than its pointwise
  form (Proposition~\ref{prop:strict}).
\end{enumerate}

The paper is organised as follows. Section~\ref{sec:main-assumptions} sets out the model
and its assumptions. Section~\ref{sec:main} states the main results.
Section~\ref{sec:wp-proof} carries out the continuation proof.
Section~\ref{sec:examples} develops the examples. Section~\ref{sec:conclusion} concludes
and identifies open extensions.
Appendix~\ref{app:lawflow} constructs the Borel law flow used throughout.

